\chapter{Math and Core Geometry Types}
\label{ch:math-core-types}

Every CNA class covered later in this book --- \cnaclass{SpriteBatch}, \cnaclass{Model},
every stock effect, every backend --- passes values of the types in this chapter across its
API. They live directly under \texttt{include/Microsoft/Xna/Framework/} (not in the
\texttt{Graphics} sub-namespace), are all value types (C\texttt{++} \texttt{struct}s, matching
XNA's own \texttt{struct} semantics rather than reference-type classes, with the sole
exception of \cnaclass{BoundingFrustum}, discussed below), and each carries a private
\texttt{getDebugDisplayStringProperty()} and \texttt{CheckForNaNs()} helper ported from FNA's
own debug-visualizer tooling.

\section{Vectors: Vector2, Vector3, Vector4}

The three vector types share one shape: public fields (\texttt{X}, \texttt{Y}, and as
applicable \texttt{Z}, \texttt{W}), a set of static constants (\texttt{Zero}, \texttt{One},
\texttt{UnitX}/\texttt{UnitY}/\ldots), a default zero-initializing constructor, a per-component
constructor, and an \texttt{explicit} single-scalar broadcast constructor
(\texttt{Vector2(float value)} sets both components to \texttt{value}). \cnaclass{Vector3} and
\cnaclass{Vector4} additionally provide composing constructors from lower-dimensional vectors
plus extra scalars --- \texttt{Vector3(Vector2, float z)}, \texttt{Vector4(Vector3, float w)},
\texttt{Vector4(Vector2, float z, float w)} --- letting code build a higher-dimensional
vector out of a lower one without manually unpacking components.

\cnaclass{Vector3} alone defines direction constants --- \texttt{Up}, \texttt{Down},
\texttt{Right}, \texttt{Left}, \texttt{Forward}, \texttt{Backward} --- with an explicit doc
comment that \texttt{Forward}/\texttt{Backward} assume a \emph{right-handed} coordinate
system (\texttt{Forward = (0,0,-1)}, \texttt{Backward = (0,0,1)}), a detail worth remembering
whenever this book's later chapters discuss camera or model transforms.

Every static method that can meaningfully be either ``return a new value'' or ``write into an
existing one'' is provided \emph{both} ways --- a value-returning overload and an output-ref
overload writing into a \texttt{result\&} parameter --- mirroring XNA's own C\# \texttt{ref}
method-overload convention exactly: \texttt{Add}, \texttt{Clamp}, \texttt{Distance}/%
\texttt{DistanceSquared}, \texttt{Divide}, \texttt{Dot}, \texttt{Lerp}, \texttt{Max}/%
\texttt{Min}, \texttt{Multiply}, \texttt{Negate}, \texttt{Normalize}, \texttt{Reflect},
\texttt{SmoothStep}, \texttt{Subtract}, and the curve-interpolation pair
\texttt{Barycentric}/\texttt{CatmullRom}/\texttt{Hermite}. \texttt{Transform} additionally
comes in by-\cnaclass{Matrix} and by-\cnaclass{Quaternion} forms, each with single-value and
full-array/range overloads; \texttt{TransformNormal} (Vector2/Vector3 only) applies a matrix
transform with translation stripped out, the form correct for transforming surface normals
rather than positions. Operators (\texttt{==}, \texttt{!=}, unary/binary \texttt{-},
component-wise and scalar \texttt{*}/\texttt{/}, and \texttt{+}) are implemented as
\texttt{friend} functions; \cnaclass{Vector3} alone adds \texttt{NOXNA} in-place
\texttt{operator+=}/\texttt{operator-=} convenience overloads beyond the strict XNA surface.

Both \cnaclass{Matrix} and \cnaclass{Quaternion} are only forward-declared in these headers
rather than included, specifically to avoid a circular include --- \texttt{Matrix.hpp} and
\texttt{Quaternion.hpp} both depend on the vector types in the other direction.

\section{Matrix: a 4x4, and one CNA-specific bridge method}

\cnaclass{Matrix} stores sixteen public \texttt{float} fields named \texttt{M11}
through \texttt{M44} (row-major naming: \texttt{M<row><col>}), with direction-vector
convenience accessors --- \texttt{getBackwardProperty}/\texttt{setBackwardProperty} (the
third row), \texttt{getForwardProperty} (negated third row), \texttt{getUpProperty} (second
row), \texttt{getDownProperty} (negated second row), \texttt{getLeftProperty}/%
\texttt{getRightProperty} (negated/plain first row), and
\texttt{getTranslationProperty}/\texttt{setTranslationProperty} (the fourth row). Unlike the
vector types' \texttt{Zero}/\texttt{One} fields, \texttt{Identity} is exposed as a static
\emph{function}, \texttt{getIdentityProperty()}, not a field constant.

The static factory surface is the largest of any type in this chapter: billboard and
constrained-billboard construction, axis-angle and quaternion and yaw/pitch/roll rotation,
look-at, orthographic (centered and off-center), perspective (plain, field-of-view, and
off-center), per-axis rotation (\texttt{CreateRotationX/Y/Z}), scale (scalar, per-axis, and
\cnaclass{Vector3} overloads), shadow, translation (scalar and \cnaclass{Vector3}
overloads), reflection, and world-matrix construction --- each, again, with a value-returning
and an output-ref overload. Arithmetic statics (\texttt{Add}, \texttt{Divide},
\texttt{Invert}, \texttt{Lerp}, \texttt{Multiply}, \texttt{Negate}, \texttt{Subtract},
\texttt{Transpose}) and the instance method \texttt{Decompose(scale\&, rotation\&,
translation\&) -> bool} round out the surface.

One method exists purely to bridge XNA's conceptual row-major matrix into the graphics
backends covered in Part~IV of this volume: \texttt{NOXNA void ToColumnMajor(float
out[16]) const}, documented as returning ``this matrix in the transposed column-major form
expected by OpenGL-style uniform uploads.'' OpenGL and Vulkan-style shader uniform layouts
expect column-major data; without this bridge, every backend's \cnaclass{Effect}/shader
implementation would need to re-derive the same transpose independently. Centralizing it on
\cnaclass{Matrix} itself is a small but telling example of where CNA allows itself exactly
one, clearly-marked graphics-motivated extension on an otherwise pure math type, rather than
scattering the same transpose logic across eleven backend implementations.

\section{Quaternion, Plane, and Ray}

\cnaclass{Quaternion} follows the same value/output-ref-pair convention as the vectors:
\texttt{Add}, \texttt{Concatenate} (value1 followed by value2 --- order matters for
quaternion composition), \texttt{Conjugate}, \texttt{CreateFromAxisAngle}/%
\texttt{CreateFromRotationMatrix}/\texttt{CreateFromYawPitchRoll}, \texttt{Divide},
\texttt{Dot}, \texttt{Inverse}, \texttt{Lerp} (normalized), \texttt{Slerp},
\texttt{Multiply} (quaternion-by-quaternion and quaternion-by-scalar), \texttt{Negate}, and
\texttt{Normalize}, plus its single static constant, \texttt{Identity}.

\cnaclass{Plane} stores a \texttt{Vector3 Normal} and a \texttt{float D} (distance from
origin), with constructors from a \cnaclass{Vector4} (XYZ as normal, W as distance), from an
explicit normal-plus-distance pair, from three points, or from four raw scalars.
\texttt{DotCoordinate}/\texttt{DotNormal} implement the two distinct dot-product forms a
plane test needs; \texttt{Intersects} overloads against \cnaclass{BoundingBox},
\cnaclass{BoundingSphere}, and \cnaclass{BoundingFrustum} all return
\cnaclass{PlaneIntersectionType} (\S\ref{sec:bounding-volumes}). \cnaclass{Plane} declares
\texttt{friend class BoundingFrustum}, because frustum plane extraction
(\S\ref{sec:bounding-volumes}) needs to construct and normalize planes directly rather than
only through \cnaclass{Plane}'s public surface.

\cnaclass{Ray} is the one type in this chapter whose XNA-to-C\texttt{++} translation is more
idiomatic than literal: XNA's \texttt{Intersects} methods return a nullable
\texttt{float?} in C\#; CNA's equivalents return \texttt{std::optional<float>} directly ---
the intersection distance if one exists, \texttt{std::nullopt} otherwise --- against
\cnaclass{BoundingBox}, \cnaclass{BoundingSphere}, \cnaclass{Plane}, and
\cnaclass{BoundingFrustum}, each with a by-ref \texttt{std::optional<float>\&} output-parameter
twin as well.

\section{Rectangle and Point}

\cnaclass{Rectangle} stores \texttt{intcs X, Y, Width, Height} --- the
\texttt{SharpRuntime} \texttt{intcs} alias from Chapter~\ref{ch:design-philosophy}'s type-alias
table, not a raw C\texttt{++} \texttt{int} --- with derived get-only properties
\texttt{Left}/\texttt{Right}/\texttt{Top}/\texttt{Bottom}, a get/set \texttt{Location}
(\cnaclass{Point}), and a get-only \texttt{Center} (\cnaclass{Point}, ``rounded down for odd
sizes''). \texttt{Contains} and \texttt{Intersects} each come in value-returning and
output-ref forms against a coordinate pair, a \cnaclass{Point}, or another
\cnaclass{Rectangle}; \texttt{Offset} and \texttt{Inflate} mutate in place; static
\texttt{Intersect}/\texttt{Union} compute a new rectangle from two inputs.

\cnaclass{Point} is the integer counterpart to \cnaclass{Vector2} --- \texttt{intcs X, Y} ---
with component-wise \texttt{+}, \texttt{-}, \texttt{*}, and \texttt{/} operators implemented
as friend functions. Note for anyone diffing against a stricter reading of the real XNA~4.0
\cnaclass{Point} surface: CNA's component-wise \texttt{Point * Point} and \texttt{Point /
Point} operators are present here regardless of whether the original XNA assembly's public
operator surface defines them identically --- a detail worth a direct \texttt{xna4-spec}
diff (Volume~II's \texttt{xna4-spec} auditing chapter) rather than an assumption, if exact
operator-surface parity matters for your use case.

\section{Color: 140 named constants and one packed layout}

\cnaclass{Color} implements CNA's \texttt{Graphics::PackedVector::IPackedVectorT<UInt32>}
and \texttt{System::IEquatable<Color>}, storing its four \texttt{bytecs} channels
(\texttt{R}, \texttt{G}, \texttt{B}, \texttt{A}) packed into a single \texttt{PackedValue}
(\texttt{UInt32}) in \texttt{AABBGGRR} byte order --- a layout detail
Chapter~\ref{ch:design-philosophy} names explicitly as one of the ``behavior fidelity''
details CNA is required to preserve exactly, since backend code and any raw-bytes
interoperability depends on it.

Roughly 140 named XNA/X11 web-color static constants are defined --- \texttt{Transparent},
\texttt{AliceBlue}, ..., \texttt{YellowGreen}, including \texttt{CornflowerBlue}, the
traditional XNA project-template clear color already seen in
Chapter~\ref{ch:first-game} --- each documented with its exact RGBA byte values. Construction
accepts a \cnaclass{Vector4} (clamped to $[0,1]$, scaled to $[0,255]$), a \cnaclass{Vector3}
(alpha forced to 255), three or four \texttt{float}s in $[0,1]$, or three or four
\texttt{intcs} in $[0,255]$ (clamped) --- plus two \texttt{NOXNA}-marked byte-based
constructors explicitly documented as ``not part of the XNA 4.0 API. CNA convenience
overload.'' \texttt{Lerp}, \texttt{FromNonPremultiplied}, and \texttt{Multiply(Color,
float)} round out the static surface; a \texttt{NOXNA} commutative \texttt{float * Color}
operator is added alongside the standard \texttt{Color * float}, again explicitly documented
as a CNA-only convenience not present in XNA/FNA.

\section{MathHelper: constants and helpers, some public only because C\texttt{++} has no assembly visibility}

\cnaclass{MathHelper} is static-only (its default constructor is \texttt{= delete}, tagged
\texttt{NOXNA}), providing \texttt{Pi}/\texttt{PiOver2}/\texttt{PiOver4}/\texttt{TwoPi}/%
\texttt{E}/\texttt{Log10E}/\texttt{Log2E} as \texttt{constexpr float} constants, plus
\texttt{MachineEpsilonFloat} computed once via a private \texttt{GetMachineEpsilonFloat()}.
Its own header comment is unusually direct about a translation artifact: ``FNA marks the
four members below as `internal'; C++ has no assembly-scope visibility, so they are public
here. Other framework types (e.g.\ Vector2, Matrix) rely on them.'' Those four --- an integer
\texttt{Clamp} overload, \texttt{WithinEpsilon}, \texttt{ClosestMSAAPower}, and
\texttt{MachineEpsilonFloat} itself --- are effectively framework-internal helpers that
simply have no way to be hidden from a caller in C\texttt{++} the way FNA hides them from
other assemblies in C\#. \texttt{ClosestMSAAPower(intcs)} in particular is worth flagging as
a graphics-relevant helper hiding in a math-only-looking class: it returns the closest valid
MSAA sample-count power-of-two at or below the requested value, mirroring FNA's own internal
MSAA helper, and is used wherever a backend needs to clamp a requested multisample count to
something the underlying graphics API actually supports.

\section{Bounding volumes: Box, Sphere, Frustum}
\label{sec:bounding-volumes}

\cnaclass{BoundingBox}, \cnaclass{BoundingSphere}, and \cnaclass{BoundingFrustum} share a
consistent overload pattern: \texttt{Contains} against every other bounding-volume type and
against a raw \cnaclass{Vector3} returns \cnaclass{ContainmentType} (\texttt{Disjoint},
\texttt{Contains}, or \texttt{Intersects}); \texttt{Intersects} against a \cnaclass{Plane}
returns \cnaclass{PlaneIntersectionType} (\texttt{Front}, \texttt{Back}, or
\texttt{Intersecting} --- the doc comments clarify that \texttt{Front}/\texttt{Back} mean
the volume lies entirely within one half-space with no actual crossing, while
\texttt{Intersecting} means the plane genuinely cuts through the volume); and
\texttt{Intersects} against a \cnaclass{Ray} returns \texttt{std::optional<float>}. Both
value-returning and output-ref forms exist throughout, matching the vector/matrix convention.

\cnaclass{BoundingBox} stores \texttt{Vector3 Min, Max} and a \texttt{CornerCount = 8}
constant; \texttt{GetCorners()} has both a value-returning form and a fill-existing-vector
form (documented as requiring the vector to already hold at least \texttt{CornerCount}
elements). \cnaclass{CHECKLIST.md}'s own known-deviations table flags one project-wide
inconsistency that happens to touch this class: \texttt{BoundingBox.cpp} throws
\texttt{System::ArgumentOutOfRangeException} directly for an out-of-range indexer access, the
same choice \texttt{VertexBuffer.cpp} and \texttt{NetworkSessionProperties.cpp} make ---
whereas \texttt{Input::Touch::TouchCollection} deliberately throws \texttt{std::out\_of\_range}
instead for the analogous case. The checklist names this explicitly as a known, unresolved,
project-wide inconsistency rather than a bug specific to any one file --- worth knowing if
your own code catches a specific exception type at one of these call sites.

\cnaclass{BoundingSphere} stores \texttt{Vector3 Center} and \texttt{float Radius}, with a
\texttt{Transform(Matrix)} method whose doc comment is precise about a real geometric
limitation: it applies only translation and (uniform) scale from the given matrix, matching
XNA's own semantics --- a sphere transformed by a shearing or non-uniform-scale matrix does
not become an ellipsoid, because \cnaclass{BoundingSphere} has no representation for one.

\cnaclass{BoundingFrustum} is the one bounding-volume type implemented as a C\texttt{++}
\texttt{class} rather than a \texttt{struct}, because unlike the other two it owns real
cached derived state: constructed from a single combined view-projection \cnaclass{Matrix},
it privately caches an \texttt{std::array<Vector3, 8>} of corners and an
\texttt{std::array<Plane, 6>} of the six extracted clip planes (\texttt{Near}, \texttt{Far},
\texttt{Left}, \texttt{Right}, \texttt{Top}, \texttt{Bottom}, each exposed as a get-only
\cnaclass{Plane} property), recomputed whenever its \texttt{Matrix} property is reassigned.
Private helpers \texttt{CreateCorners()}, \texttt{CreatePlanes()}, a static
\texttt{NormalizePlane(Plane\&)}, and a static \texttt{IntersectionPoint(...)} (finding a
frustum corner as the intersection of three planes) implement this derivation; equality
(\texttt{operator==}/\texttt{!=}) is defined purely in terms of the underlying \texttt{Matrix},
not the cached corners/planes, since the latter are wholly determined by the former.

\section{Curves: Curve, CurveKey, and their supporting enums}

\cnaclass{Curve} holds an ordered \cnaclass{CurveKeyCollection} of \cnaclass{CurveKey} points
and evaluates the resulting piecewise curve at an arbitrary position via
\texttt{Evaluate(float) -> float}; \texttt{IsConstant} reports whether the curve has zero or
one keys (in which case evaluation is trivially constant). Two independently-settable
\cnaclass{CurveLoopType} properties, \texttt{PreLoop} and \texttt{PostLoop}, govern what
happens when \texttt{Evaluate} is asked for a position before the first key or after the
last one: \texttt{Constant} clamps to the boundary key's value, \texttt{Linear} extrapolates
linearly, \texttt{Cycle} wraps the curve's own span, \texttt{CycleOffset} wraps \emph{and}
adds the accumulated first/last-key value difference per cycle (useful for a monotonically
increasing quantity like distance traveled along a looping path), and \texttt{Oscillate}
ping-pongs direction each cycle. A private \texttt{GetCurvePosition(float) -> float},
together with \texttt{GetNumberOfCycle(float) -> int}, implements this remapping before the
underlying piecewise evaluation runs.

Tangent handling is exposed at three granularities: \texttt{ComputeTangents(CurveTangent)}
sets both in and out tangents for every key to the same \cnaclass{CurveTangent} mode;
\texttt{ComputeTangents(CurveTangent in, CurveTangent out)} sets them independently for every
key; and \texttt{ComputeTangent(keyIndex, ...)} (both single-mode and in/out-mode overloads)
recomputes tangents for one specific key. The three \cnaclass{CurveTangent} modes are
\texttt{Flat} (tangent forced to zero), \texttt{Linear} (tangent equal to the difference to
the neighboring key's value), and \texttt{Smooth} (derived from both neighboring keys).

Each \cnaclass{CurveKey} carries an immutable \texttt{Position} (set only at construction),
a mutable \texttt{Value}, mutable \texttt{TangentIn}/\texttt{TangentOut}, and a
\cnaclass{CurveContinuity} (\texttt{Smooth} --- interpolate normally between this key and the
next --- or \texttt{Step}, where any position between this key and the next simply returns
this key's value, no interpolation). \cnaclass{CurveKeyCollection} keeps its keys sorted by
position at all times --- its \texttt{Add} method is documented as ``adds a key while
preserving ascending position order,'' i.e.\ an insertion-sorted add, not an append --- and
adds the same \texttt{NOXNA} C\texttt{++}-iterator convenience surface
(\texttt{iterator}/\texttt{const\_iterator}/\texttt{begin}/\texttt{end}/\texttt{cbegin}/%
\texttt{cend}) already seen on \cnaclass{GameComponentCollection} in
Chapter~\ref{ch:game-loop}.

\section{A note on scope: this chapter versus NEXT.md}

It is worth being explicit about what this chapter does \emph{not} claim. CNA's root-level
\texttt{NEXT.md} --- described in Chapter~\ref{ch:what-is-cna} as the single most reliable,
actively-maintained bug list in the repository --- is almost entirely occupied by the Audio
subsystem (XACT/FAudio parity work) and by graphics-backend porting phases; where it mentions
\texttt{Matrix}, \texttt{Vector2}/\texttt{3}, or \texttt{Plane} by name, it is nearly always in
the context of a \emph{backend's} shader or transform implementation (for example, a D3D9
matrix-transform bug affecting \cnaclass{SpriteBatch} layer-depth sorting, covered in
Chapter~\ref{ch:direct3d-backends}), not the core \cnaclass{Microsoft::Xna::Framework} types
covered in this chapter. As of the source material this book draws from, none of
\cnaclass{GameTime}, \cnaclass{Curve}/\cnaclass{CurveKey}, \cnaclass{GameComponent}, or the
math/bounding-volume types themselves carry any open, named bug in \texttt{NEXT.md} --- their
stability is precisely why this chapter can describe them as settled reference material rather
than work in progress.
